\documentclass[10pt]{beamer}
\usetheme{Madrid}
\usecolortheme{default}
\usepackage{graphicx}
\usepackage{booktabs}
\usepackage{multirow}
\usepackage{xcolor}

\title[MINDSET Employment Methods]{Methodological Development: \\Employment-Only Mode for MINDSET}
\subtitle{Custom Scripts and Workflow for 469 Strategy Analysis}
\author{Felipe Esteves}
\institute{RAND Corporation}
\date{\today}

\begin{document}

\frame{\titlepage}

% ============================================
% SECTION 1: STARTING POINT
% ============================================

\begin{frame}{Starting Point: Base MINDSET Model}
\textbf{What We Had:}

\vspace{0.3cm}

\begin{itemize}
    \item Base MINDSET model (World Bank)
    \item Full dynamic CGE model with multiple modules
    \item Designed for: energy taxes, trade policy, full equilibrium analysis
    \item Runtime: 10-30 minutes per scenario
    \item \textcolor{red}{\textbf{Problem:}} Too slow for 469 scenarios (78+ hours)
\end{itemize}

\vspace{0.5cm}

\textbf{Our Goal:}

\vspace{0.3cm}

\begin{itemize}
    \item Estimate \textbf{employment impacts only} (not full CGE)
    \item Process 469 investment strategy files
    \item Output: Direct vs indirect jobs by ISIC sector
    \item Target runtime: $<$1 hour total
\end{itemize}

\vspace{0.5cm}

\textbf{Solution:} Build custom employment-only batch processing workflow
\end{frame}

% ============================================
% SECTION 2: KEY CHALLENGES
% ============================================

\begin{frame}{Three Key Challenges We Solved}

\textbf{Challenge 1: Strategy File Format Mismatch}
\begin{itemize}
    \item Strategy files use 26 Eora26 sectors
    \item MINDSET uses 120 GLORIA products
    \item \textcolor{blue}{\textbf{Solution:}} Built product allocation system with crosswalk
\end{itemize}

\vspace{0.4cm}

\textbf{Challenge 2: Investment Interpretation}
\begin{itemize}
    \item Ambiguous: Is investment "by" sector or "purchasing" products?
    \item Wrong interpretation = circular logic
    \item \textcolor{blue}{\textbf{Solution:}} Implemented "strategy as purchases" logic
\end{itemize}

\vspace{0.4cm}

\textbf{Challenge 3: Output Aggregation}
\begin{itemize}
    \item Need results in 20 ISIC sectors (for dissertation)
    \item MINDSET outputs 120 products
    \item \textcolor{blue}{\textbf{Solution:}} Built two-stage aggregation pipeline
\end{itemize}

\end{frame}

% ============================================
% SECTION 3: METHODOLOGICAL INCREMENT 1
% ============================================

\begin{frame}[fragile]{Increment 1: Product Allocation System}
\textbf{Problem:} Strategy files specify investment by Eora26 sector, but MINDSET needs product-level demand

\vspace{0.3cm}

\textbf{What We Built:}

\vspace{0.2cm}

\textbf{Step 1:} Created crosswalk matrix
\begin{itemize}
    \item File: \texttt{Data/GLORIA-Eora26 - Crosswalk.xlsx}
    \item Structure: 120 products $\times$ 26 sectors (binary indicators)
    \item Example: "Wheat" → "Agriculture", "Steel" → "Manufacturing"
\end{itemize}

\vspace{0.2cm}

\textbf{Step 2:} Built allocation script
\begin{itemize}
    \item File: \texttt{sector\_to\_product\_allocation.Rmd}
    \item Logic: Divide sector value evenly among linked products
    \item Formula: $\text{Value}_{\text{product}} = \frac{\text{Value}_{\text{sector}}}{\text{\# products in sector}}$
\end{itemize}

\vspace{0.2cm}

\textbf{Step 3:} Generated country-specific distributions
\begin{itemize}
    \item Input: 7 cost structure files (\texttt{cost\_str\_*.xlsx})
    \item Output: \texttt{sector\_to\_product\_distribution\_allcountries.csv}
    \item Validation: Values sum to 1.0 per country per strategy ✓
\end{itemize}

\end{frame}

\begin{frame}[fragile]{Product Allocation: Example}
\textbf{Input (Strategy File):}

\vspace{0.2cm}

\begin{table}[h]
\centering
\small
\begin{tabular}{lr}
\toprule
\textbf{Eora26 Sector} & \textbf{Share} \\
\midrule
Agriculture & 0.35 \\
Manufacturing & 0.45 \\
Construction & 0.20 \\
\bottomrule
\end{tabular}
\end{table}

\vspace{0.3cm}

\textbf{After Allocation (GLORIA Products):}

\vspace{0.2cm}

\begin{table}[h]
\centering
\footnotesize
\begin{tabular}{lrl}
\toprule
\textbf{Product} & \textbf{Share} & \textbf{Calculation} \\
\midrule
Wheat & 0.0438 & $0.35 / 8$ products in Agriculture \\
Maize & 0.0438 & $0.35 / 8$ \\
... (6 more agri products) & ... & ... \\
Food processing & 0.0100 & $0.45 / 45$ products in Manufacturing \\
Textiles & 0.0100 & $0.45 / 45$ \\
... (43 more mfg products) & ... & ... \\
Construction services & 0.2000 & $0.20 / 1$ product in Construction \\
\bottomrule
\end{tabular}
\end{table}

\textbf{Validation:} $\sum = 1.0$ ✓
\end{frame}

\begin{frame}{Intuition: Why Product Allocation Matters}
\textbf{Economic Question:} When you invest \$1M in "Manufacturing", which products get purchased?

\vspace{0.3cm}

\textbf{Real-World Example:}

Building a factory requires:
\begin{itemize}
    \item Steel (20\% of cost)
    \item Cement (15\%)
    \item Machinery (30\%)
    \item Electrical equipment (20\%)
    \item Other materials (15\%)
\end{itemize}

\vspace{0.3cm}

\textbf{Without Product Allocation:}
\begin{itemize}
    \item[\textcolor{red}{✗}] Model treats all "Manufacturing" as homogeneous
    \item[\textcolor{red}{✗}] Misses sector-specific linkages
    \item[\textcolor{red}{✗}] Cannot trace which industries benefit
\end{itemize}

\vspace{0.3cm}

\textbf{With Product Allocation:}
\begin{itemize}
    \item[\textcolor{blue}{✓}] Demand distributed across correct products
    \item[\textcolor{blue}{✓}] Captures realistic supply chains
    \item[\textcolor{blue}{✓}] Employment effects in right sectors
\end{itemize}

\end{frame}

% ============================================
% SECTION 4: METHODOLOGICAL INCREMENT 2
% ============================================

\begin{frame}[fragile]{Increment 2: "Strategy as Purchases" Logic}
\textbf{Conceptual Issue:}

Investment of \$1M in "Agriculture" could mean:
\begin{enumerate}
    \item Investment \textbf{by} agricultural sector → creates agri output
    \item Investment \textbf{purchases} agricultural products → creates demand
\end{enumerate}

\vspace{0.3cm}

\textbf{Why This Matters:}

\vspace{0.2cm}

\textcolor{red}{\textbf{Wrong (Option 1):}} Circular logic - investment creates its own supply
\begin{itemize}
    \item Violates input-output accounting
    \item Double-counts employment
\end{itemize}

\vspace{0.2cm}

\textcolor{blue}{\textbf{Correct (Option 2):}} Investment creates final demand
\begin{itemize}
    \item Follows standard IO theory
    \item Leontief model traces supply chain effects
    \item Separates direct from indirect employment
\end{itemize}

\end{frame}

\begin{frame}[fragile]{Implementation: Investment Converter}
\textbf{What We Modified:}

\vspace{0.3cm}

\textbf{File:} \texttt{RunMINDSET\_EmploymentOnly\_BATCH\_FINAL.py}, lines 126-145

\vspace{0.2cm}

\begin{verbatim}
# STEP 2: Investment Converter
# (Sector Investment → Product Demand)

INV_CONV = Scenario.set_inv_conv_adj(MRIO_BASE)
INV_model = invest(MRIO_BASE, INV_CONV, Scenario)
INV_model.build_inv_output_elas()
INV_model.calc_inv_share()
INV_model.calc_dy_inv_exog(Scenario.inv_exog)

# Convert to vector: INITIAL FINAL DEMAND
dy_inv_exog = MRIO_df_to_vec(
    INV_model.dy_inv_exog, ...
)
\end{verbatim}

\vspace{0.3cm}

\textbf{Key Point:} Uses MINDSET's \texttt{investment.py} module to map investment → product demand via \texttt{INV\_CONV} matrix

\end{frame}

\begin{frame}{Intuition: Why "Purchases" vs "Production" Matters}
\textbf{Concrete Example:} \$1M investment in "Agriculture"

\vspace{0.3cm}

\textbf{Wrong Interpretation (Investment BY Sector):}

\vspace{0.2cm}

\textcolor{red}{\textbf{Scenario 1:}} Farmers invest in their own farms
\begin{itemize}
    \item Farmers buy tractors, seeds, fertilizer
    \item This creates demand for \textbf{manufacturing} (tractors), \textbf{chemicals} (fertilizer)
    \item But model would credit jobs to \textbf{agriculture only}
    \item \textcolor{red}{Result: Misses supply chain jobs}
\end{itemize}

\vspace{0.4cm}

\textbf{Correct Interpretation (Investment PURCHASES Products):}

\vspace{0.2cm}

\textcolor{blue}{\textbf{Scenario 2:}} Someone invests to buy agricultural products
\begin{itemize}
    \item Creates demand for food, crops, livestock
    \item Supply chain activates: farmers → processors → distributors
    \item Leontief model traces \textbf{all} backward linkages
    \item \textcolor{blue}{Result: Captures full employment impact}
\end{itemize}

\vspace{0.4cm}

\textbf{Bottom Line:} Investment creates \textbf{final demand}, not direct production

\end{frame}

% ============================================
% SECTION 5: METHODOLOGICAL INCREMENT 3
% ============================================

\begin{frame}{Increment 3: Employment-Only Batch Processing}
\textbf{What We Built:}

\vspace{0.3cm}

\textbf{Script:} \texttt{RunMINDSET\_EmploymentOnly\_BATCH\_FINAL.py}

\vspace{0.2cm}

\textbf{Key Modifications from Base MINDSET:}

\vspace{0.2cm}

\begin{enumerate}
    \item \textbf{Disabled unused modules}
    \begin{itemize}
        \item Energy module (not needed)
        \item Trade substitution (not needed)
        \item Price effects (not needed)
        \item Household response (not needed)
    \end{itemize}

    \vspace{0.2cm}

    \item \textbf{Kept essential modules}
    \begin{itemize}
        \item Input-Output (Leontief model)
        \item Employment (multipliers)
        \item Investment (converter)
    \end{itemize}

    \vspace{0.2cm}

    \item \textbf{Added direct vs indirect separation}
    \begin{itemize}
        \item Direct: $\Delta E_{\text{direct}} = m \circ \Delta y$
        \item Total: $\Delta E_{\text{total}} = m \circ (\mathbf{L} \times \Delta y)$
        \item Indirect: $\Delta E_{\text{indirect}} = \Delta E_{\text{total}} - \Delta E_{\text{direct}}$
    \end{itemize}
\end{enumerate}

\end{frame}

\begin{frame}[fragile]{Batch Processing Workflow}
\textbf{Pseudocode of Main Loop:}

\vspace{0.2cm}

\begin{verbatim}
# ONE-TIME INITIALIZATION (heavy computation)
MRIO_BASE = load_gloria_data()
L_BASE = calculate_leontief_inverse()  # (I-A)^-1
empl_multipliers = calculate_multipliers()

# LOOP: 469 iterations (fast)
for strategy in strategy_files:
    # Load investment
    dy = load_strategy_investment(strategy)

    # Apply multipliers (matrix operations)
    dq_total = L_BASE @ dy           # Output change
    jobs_total = empl_mult * dq_total
    jobs_direct = empl_mult * dy
    jobs_indirect = jobs_total - jobs_direct

    # Aggregate to ISIC sectors
    results = aggregate_to_ISIC(jobs_total)

    # Save
    all_results.append(results)
\end{verbatim}

\vspace{0.3cm}

\textbf{Performance:} Initialization 5 min, loop 0.1 sec/strategy $\approx$ 50 min total

\end{frame}

\begin{frame}{Intuition: Direct vs Indirect Employment}
\textbf{Example:} Government builds a \$10M highway

\vspace{0.3cm}

\textbf{Direct Employment (First-Round):}
\begin{itemize}
    \item Construction workers building the highway
    \item Engineers designing the road
    \item Heavy equipment operators
    \item \textbf{Formula:} $\Delta E_{\text{direct}} = m \circ \Delta y$ (initial demand only)
    \item \textbf{Typical:} 60-70\% of total jobs
\end{itemize}

\vspace{0.4cm}

\textbf{Indirect Employment (Supply Chain):}
\begin{itemize}
    \item \textbf{Steel mills} producing rebar for concrete
    \item \textbf{Cement plants} supplying materials
    \item \textbf{Truckers} transporting materials
    \item \textbf{Equipment manufacturers} making bulldozers
    \item \textbf{Fuel refineries} supplying diesel
    \item \textbf{Formula:} $\Delta E_{\text{indirect}} = m \circ [(\mathbf{L} - \mathbf{I}) \times \Delta y]$ (multiplier effects)
    \item \textbf{Typical:} 30-40\% additional jobs
\end{itemize}

\vspace{0.3cm}

\textbf{Key Insight:} Leontief model ($\mathbf{L}$) automatically traces all supply chain effects

\end{frame}

\begin{frame}{Intuition: How the Leontief Multiplier Works}
\textbf{Simple 3-Sector Example:}

\vspace{0.3cm}

\textbf{Scenario:} \$1M investment in Construction

\vspace{0.2cm}

\textbf{Round 1 (Direct):}
\begin{itemize}
    \item Construction sector receives \$1M demand
    \item Hires workers, buys materials
    \item Creates 15 jobs directly
\end{itemize}

\vspace{0.2cm}

\textbf{Round 2 (1st-order indirect):}
\begin{itemize}
    \item Construction needs \$300k steel, \$200k cement
    \item Steel sector increases production → hires 3 workers
    \item Cement sector increases production → hires 2 workers
\end{itemize}

\vspace{0.2cm}

\textbf{Round 3 (2nd-order indirect):}
\begin{itemize}
    \item Steel production needs \$50k iron ore
    \item Cement production needs \$30k limestone
    \item Mining sectors hire 1 more worker
\end{itemize}

\vspace{0.2cm}

\textbf{Round 4, 5, 6...} (diminishing effects)

\vspace{0.3cm}

\textbf{Leontief inverse} $\mathbf{L} = (\mathbf{I} - \mathbf{A})^{-1}$ computes \textbf{sum of infinite rounds}

Total jobs = $15 + 5 + 1 + 0.3 + ... \approx$ \textbf{21 jobs}

\end{frame}

% ============================================
% SECTION 6: METHODOLOGICAL INCREMENT 4
% ============================================

\begin{frame}{Increment 4: ISIC Sector Aggregation Pipeline}
\textbf{Challenge:} Map 120 GLORIA products → 20 ISIC Rev. 4 sections

\vspace{0.3cm}

\textbf{Two-Stage Aggregation:}

\vspace{0.2cm}

\textbf{Stage 1:} GLORIA (120) → Eora26 (26)
\begin{itemize}
    \item Used product allocation crosswalk (inverse mapping)
    \item Aggregates employment within Eora26 sectors
\end{itemize}

\vspace{0.2cm}

\textbf{Stage 2:} Eora26 (26) → ISIC Rev. 4 (20)
\begin{itemize}
    \item File: \texttt{Data/Eora26-ISIC\_crosswalk.xlsx} (created by team)
    \item Maps Eora26 to standard ISIC sections A-T
    \item Handles many-to-one relationships
\end{itemize}

\vspace{0.3cm}

\textbf{Implementation:}

\vspace{0.2cm}

Created aggregation function in batch script:
\begin{itemize}
    \item Input: 120-element employment vector
    \item Output: 20-element ISIC employment vector
    \item Validation: No jobs lost in aggregation ✓
\end{itemize}

\end{frame}

\begin{frame}[fragile]{ISIC Aggregation: Code Structure}
\textbf{File:} \texttt{RunMINDSET\_EmploymentOnly\_BATCH\_FINAL.py}, lines 231-252

\vspace{0.2cm}

\begin{verbatim}
# Calculate employment by sector
empl_by_sector = empl_results.groupby('Sec_ID').agg({
    'Direct_Jobs': 'sum',
    'Indirect_Jobs': 'sum',
    'Total_Jobs': 'sum'
}).reset_index()

# Merge with sector names
empl_by_sector = empl_by_sector.merge(
    MRIO_BASE.P[['Lfd_Nr', 'Sector_names']],
    left_on='Sec_ID',
    right_on='Lfd_Nr',
    how='left'
)

# Map to ISIC (external crosswalk applied here)
# [Additional mapping code in post-processing]
\end{verbatim}

\vspace{0.3cm}

\textbf{Note:} ISIC mapping done in post-processing to maintain flexibility

\end{frame}

\begin{frame}{Intuition: Why ISIC Aggregation Matters}
\textbf{Problem:} GLORIA has 120 products - too granular for policy analysis

\vspace{0.3cm}

\textbf{Example Products in GLORIA:}
\begin{itemize}
    \item "Growing wheat", "Growing maize", "Growing rice", "Cattle", "Fishing", ...
    \item "Food processing", "Beverages", "Tobacco", ...
    \item "Textiles", "Wearing apparel", "Leather products", ...
\end{itemize}

\vspace{0.3cm}

\textbf{For Policy Analysis, We Need:}
\begin{itemize}
    \item \textbf{ISIC Section A (Agriculture):} All farming + fishing
    \item \textbf{ISIC Section C (Manufacturing):} All manufacturing products
    \item \textbf{ISIC Section F (Construction):} Construction activities
\end{itemize}

\vspace{0.3cm}

\textbf{Benefits of ISIC Classification:}
\begin{enumerate}
    \item \textbf{International comparability} - All countries use ISIC
    \item \textbf{Policy relevance} - Governments organize data by ISIC
    \item \textbf{Tractability} - 20 sectors vs 120 products
    \item \textbf{Alignment with national accounts} - Matches official statistics
\end{enumerate}

\vspace{0.3cm}

\textbf{Aggregation preserves total:} $\sum_{j \in \text{ISIC A}} E_j = E_{\text{Agriculture}}$

\end{frame}

\begin{frame}{Intuition: Employment Multipliers}
\textbf{What is an employment multiplier?}

\vspace{0.3cm}

\textbf{Definition:} Jobs created per \$1 million of economic output

\vspace{0.3cm}

\textbf{Formula:}
\[
m_{i,j} = \text{EMPL\_COEF}_{i,j} \times \frac{E^{\text{base}}_{i,j}}{q^{\text{base}}_{i,j}}
\]

\vspace{0.3cm}

\textbf{Intuition:}
\begin{itemize}
    \item $E^{\text{base}} / q^{\text{base}}$ = Current workers per dollar of output
    \item EMPL\_COEF = Adjustment factor from GLORIA database
    \item \textbf{Result:} Country-sector specific "jobs per \$M" rate
\end{itemize}

\vspace{0.3cm}

\textbf{Why multipliers vary:}

\vspace{0.2cm}

\begin{table}[h]
\centering
\small
\begin{tabular}{lcc}
\toprule
\textbf{Sector} & \textbf{Jobs/\$1M} & \textbf{Why?} \\
\midrule
Agriculture & 50-150 & Labor-intensive, low wages \\
Manufacturing & 10-30 & Capital-intensive, automation \\
Construction & 20-40 & Mix of labor and equipment \\
Services & 30-80 & Varies widely by type \\
\bottomrule
\end{tabular}
\end{table}

\end{frame}

% ============================================
% SECTION 7: COMPLETE WORKFLOW
% ============================================

\begin{frame}{Complete Workflow: Step-by-Step Replication}
\textbf{Phase 1: Data Preparation (One-Time)}

\vspace{0.2cm}

\begin{enumerate}
    \item \textbf{Create product allocation}
    \begin{itemize}
        \item Run: \texttt{sector\_to\_product\_allocation.Rmd}
        \item Output: \texttt{sector\_to\_product\_distribution\_allcountries.csv}
    \end{itemize}

    \vspace{0.1cm}

    \item \textbf{Calculate Leontief inverse} (if not exists)
    \begin{itemize}
        \item Run: \texttt{CALCULATE\_L\_BASE\_ONCE.py}
        \item Output: \texttt{GLORIA\_db/v57/2019/SSP/GLORIA\_L\_Base\_2019.mat}
        \item Runtime: 30-180 seconds (memory intensive)
    \end{itemize}

    \vspace{0.1cm}

    \item \textbf{Prepare ISIC crosswalk}
    \begin{itemize}
        \item File: \texttt{Data/Eora26-ISIC\_crosswalk.xlsx}
        \item Manual mapping of 26 Eora26 to 20 ISIC sections
    \end{itemize}
\end{enumerate}

\end{frame}

\begin{frame}{Complete Workflow (continued)}
\textbf{Phase 2: Batch Processing}

\vspace{0.2cm}

\begin{enumerate}
    \setcounter{enumi}{3}

    \item \textbf{Run employment batch}
    \begin{itemize}
        \item Script: \texttt{RunMINDSET\_EmploymentOnly\_BATCH\_FINAL.py}
        \item Command: \texttt{python RunMINDSET\_EmploymentOnly\_BATCH\_FINAL.py}
        \item Runtime: $\sim$50 minutes for 469 scenarios
    \end{itemize}

    \vspace{0.1cm}

    \item \textbf{Outputs generated}
    \begin{itemize}
        \item \texttt{ALL\_RESULTS\_Employment\_by\_Region.csv}
        \item \texttt{ALL\_RESULTS\_Employment\_by\_Sector.csv}
        \item \texttt{BATCH\_SUMMARY.csv}
    \end{itemize}
\end{enumerate}

\vspace{0.3cm}

\textbf{Phase 3: Post-Processing}

\vspace{0.2cm}

\begin{enumerate}
    \setcounter{enumi}{5}

    \item \textbf{Aggregate to ISIC}
    \begin{itemize}
        \item Apply Eora26-ISIC crosswalk
        \item Calculate sectoral shares
    \end{itemize}

    \item \textbf{Final consolidation}
    \begin{itemize}
        \item Merge with strategy attributes
        \item Add country names
        \item Output: \texttt{employment\_consolidated.csv}
    \end{itemize}
\end{enumerate}

\end{frame}

% ============================================
% SECTION 8: KEY FILES CREATED
% ============================================

\begin{frame}[fragile]{Key Files We Created}
\textbf{Custom Scripts (New):}

\vspace{0.2cm}

\begin{enumerate}
    \item \texttt{sector\_to\_product\_allocation.Rmd}
    \begin{itemize}
        \item Maps Eora26 sectors → GLORIA products
        \item Uses crosswalk matrix
    \end{itemize}

    \item \texttt{RunMINDSET\_EmploymentOnly\_BATCH\_FINAL.py}
    \begin{itemize}
        \item Main batch processing script
        \item 483 lines, implements employment-only mode
    \end{itemize}

    \item \texttt{CALCULATE\_L\_BASE\_ONCE.py}
    \begin{itemize}
        \item One-time Leontief inverse calculation
        \item Saves to .mat file for repeated use
    \end{itemize}
\end{enumerate}

\vspace{0.3cm}

\textbf{Data Files Created:}

\vspace{0.2cm}

\begin{enumerate}
    \setcounter{enumi}{3}

    \item \texttt{Data/GLORIA-Eora26 - Crosswalk.xlsx}
    \begin{itemize}
        \item 120 products $\times$ 26 sectors
    \end{itemize}

    \item \texttt{Data/Eora26-ISIC\_crosswalk.xlsx}
    \begin{itemize}
        \item 26 Eora26 sectors $\times$ 20 ISIC sections
    \end{itemize}

    \item \texttt{sector\_to\_product\_distribution\_allcountries.csv}
    \begin{itemize}
        \item Country-specific product allocations
    \end{itemize}
\end{enumerate}

\end{frame}

% ============================================
% SECTION 9: MODIFICATIONS TO BASE MINDSET
% ============================================

\begin{frame}{Modifications to Base MINDSET Modules}
\textbf{Modified Existing Files:}

\vspace{0.3cm}

\begin{enumerate}
    \item \texttt{SourceCode/InputOutput.py}, line 210
    \begin{itemize}
        \item \textbf{Bug fix:} Use \texttt{actual\_dims} instead of \texttt{DIMS}
        \item Issue: Mismatch between theoretical (19,680) vs actual (960) dimensions
        \item Impact: Prevents matrix dimension errors
    \end{itemize}

    \vspace{0.3cm}

    \item \textbf{No other core modules modified}
    \begin{itemize}
        \item \texttt{employment.py} - Used as-is ✓
        \item \texttt{investment.py} - Used as-is ✓
        \item \texttt{scenario.py} - Used as-is ✓
    \end{itemize}
\end{enumerate}

\vspace{0.5cm}

\textbf{Design Principle:} Minimal modifications to base MINDSET
\begin{itemize}
    \item Easier to maintain
    \item Can incorporate future MINDSET updates
    \item Custom logic isolated in batch script
\end{itemize}

\end{frame}

% ============================================
% SECTION 10: DATA QUALITY ISSUE
% ============================================

\begin{frame}{Data Quality Issue Discovered: Uganda \& Libya}
\textbf{Problem:} Employment estimates 20-100$\times$ higher than comparable countries

\vspace{0.3cm}

\begin{table}[h]
\centering
\footnotesize
\begin{tabular}{lcc}
\toprule
\textbf{Country} & \textbf{Jobs per \$1M Output} & \textbf{vs. Normal} \\
\midrule
Bulgaria, Mexico, Morocco & 10 -- 50 & Normal \\
\midrule
\textcolor{red}{\textbf{Libya}} & \textcolor{red}{\textbf{100 -- 170}} & \textcolor{red}{\textbf{5-10$\times$ higher}} \\
\textcolor{red}{\textbf{Uganda}} & \textcolor{red}{\textbf{272 -- 2,051}} & \textcolor{red}{\textbf{20-100$\times$ higher}} \\
\bottomrule
\end{tabular}
\end{table}

\vspace{0.3cm}

\textbf{Root Cause:} GLORIA v57 data quality issue
\begin{itemize}
    \item Small baseline output ($q^{\text{base}}$) values for UGA/LBY
    \item Employment multiplier formula: $m = \text{EMPL\_COEF} \times \frac{E^{\text{base}}}{q^{\text{base}}}$
    \item Small denominator → inflated multiplier
\end{itemize}

\vspace{0.3cm}

\textbf{Solution Implemented:}
\begin{itemize}
    \item Documented anomaly in methodology
    \item Flagged UGA/LBY results for sensitivity analysis
    \item Recommended: Use proxy countries (Kenya for UGA, Tunisia for LBY)
\end{itemize}

\end{frame}

% ============================================
% SECTION 11: VALIDATION CHECKS
% ============================================

\begin{frame}{Validation Checks Built Into Workflow}
\textbf{Automated Checks in Batch Script:}

\vspace{0.3cm}

\begin{enumerate}
    \item \textbf{Sum-to-one validation}
    \begin{itemize}
        \item Product allocations sum to 1.0 per strategy ✓
        \item Sectoral shares sum to 1.0 per strategy ✓
    \end{itemize}

    \vspace{0.2cm}

    \item \textbf{Accounting identity}
    \begin{itemize}
        \item Total Employment = Direct + Indirect ✓
        \item No negative employment values ✓
    \end{itemize}

    \vspace{0.2cm}

    \item \textbf{Coverage checks}
    \begin{itemize}
        \item All 469 strategies processed ✓
        \item All 7 countries present ✓
        \item All 20 ISIC sectors per strategy (including zeros) ✓
    \end{itemize}

    \vspace{0.2cm}

    \item \textbf{Reasonableness checks}
    \begin{itemize}
        \item Employment multipliers 10-100 jobs per \$1M (flagged UGA/LBY)
        \item Sectoral patterns match investment type
    \end{itemize}
\end{enumerate}

\end{frame}

% ============================================
% SECTION 12: PERFORMANCE OPTIMIZATIONS
% ============================================

\begin{frame}{Performance Optimizations Implemented}
\textbf{Key Design Decisions:}

\vspace{0.3cm}

\begin{enumerate}
    \item \textbf{One-time initialization}
    \begin{itemize}
        \item Load MRIO data once (not 469 times)
        \item Calculate Leontief inverse once
        \item Calculate employment multipliers once
        \item \textbf{Savings:} 95\% reduction in computation time
    \end{itemize}

    \vspace{0.3cm}

    \item \textbf{Vectorized operations}
    \begin{itemize}
        \item Use NumPy matrix operations (not loops)
        \item Element-wise multiplication for employment
        \item \textbf{Savings:} 100$\times$ faster than iterative
    \end{itemize}

    \vspace{0.3cm}

    \item \textbf{Sparse matrix storage} (built into base MINDSET)
    \begin{itemize}
        \item Store only non-zero MRIO elements
        \item 921,600 elements → 50,000 non-zero
        \item \textbf{Savings:} 95\% memory reduction
    \end{itemize}
\end{enumerate}

\vspace{0.3cm}

\textbf{Result:} 78 hours → 50 minutes (95$\times$ speedup)

\end{frame}

\begin{frame}{Intuition: Why Performance Optimizations Matter}
\textbf{The Computational Challenge:}

\vspace{0.3cm}

Without optimization, each scenario requires:
\begin{enumerate}
    \item Load GLORIA data from disk: 30 seconds
    \item Calculate Leontief inverse: 2-5 minutes
    \item Calculate employment multipliers: 10 seconds
    \item Run IO model: 5 seconds
    \item \textbf{Total: $\sim$10 minutes per scenario}
\end{enumerate}

\vspace{0.2cm}

$10 \text{ min} \times 469 \text{ scenarios} = \textbf{78 hours}$ 😱

\vspace{0.4cm}

\textbf{Key Insight:} Most of this is \textbf{redundant}
\begin{itemize}
    \item GLORIA data is the same for all scenarios
    \item Leontief inverse $\mathbf{L}$ is constant (same economic structure)
    \item Only $\Delta y$ (investment vector) changes per scenario
\end{itemize}

\vspace{0.4cm}

\textbf{With One-Time Initialization:}
\begin{enumerate}
    \item Load \& calculate once: 5 minutes (one-time)
    \item Per scenario: $\mathbf{L} \times \Delta y$ + $m \circ \Delta q$ = 6 seconds
    \item \textbf{Total: $5 \text{ min} + (6 \text{ sec} \times 469) = 52$ minutes} ✓
\end{enumerate}

\vspace{0.3cm}

\textbf{Lesson:} Identify constant vs variable components → compute once, reuse often

\end{frame}

% ============================================
% SECTION 13: OUTPUT STRUCTURE
% ============================================

\begin{frame}[fragile]{Output Dataset Structure}
\textbf{Final File:} \texttt{employment\_consolidated.csv}

\vspace{0.3cm}

\textbf{Dimensions:}
\begin{itemize}
    \item Rows: 9,380 (469 strategies $\times$ 20 ISIC sectors)
    \item Columns: 10
\end{itemize}

\vspace{0.3cm}

\textbf{Column Definitions:}

\vspace{0.2cm}

\begin{table}[h]
\centering
\footnotesize
\begin{tabular}{lll}
\toprule
\textbf{Column} & \textbf{Type} & \textbf{Description} \\
\midrule
country\_name & String & Full country name \\
country\_ISO & String & 3-letter ISO code \\
strategy\_id & String & e.g., "Strategy\_1004\_EGY" \\
strategy\_name & String & From attribute file \\
sector\_code & String & ISIC section (A-T) \\
sector\_name & String & ISIC description \\
direct\_jobs & Float & Direct employment \\
indirect\_jobs & Float & Indirect employment \\
total\_jobs & Float & Total = Direct + Indirect \\
share\_of\_total\_jobs & Float & Sector share (sums to 1.0) \\
\bottomrule
\end{tabular}
\end{table}

\end{frame}

% ============================================
% SECTION 14: REPLICATION CHECKLIST
% ============================================

\begin{frame}{Replication Checklist}
\textbf{To Replicate This Analysis:}

\vspace{0.3cm}

\textbf{Prerequisites:}
\begin{enumerate}
    \item GLORIA v57 database in \texttt{GLORIA\_db/}
    \item Python 3.8+ with packages: numpy, pandas, scipy, openpyxl
    \item R 4.0+ with: readxl, dplyr, tidyr (for product allocation)
    \item 469 strategy files in \texttt{GLORIA\_template/Scenarios/}
\end{enumerate}

\vspace{0.3cm}

\textbf{Step-by-Step:}
\begin{enumerate}
    \item Run \texttt{sector\_to\_product\_allocation.Rmd}
    \item Run \texttt{CALCULATE\_L\_BASE\_ONCE.py} (if L\_BASE missing)
    \item Run \texttt{RunMINDSET\_EmploymentOnly\_BATCH\_FINAL.py}
    \item Apply ISIC aggregation (post-processing)
    \item Merge with strategy attributes
\end{enumerate}

\vspace{0.3cm}

\textbf{Expected Runtime:}
\begin{itemize}
    \item Step 1: 2-5 minutes
    \item Step 2: 30-180 seconds (one-time)
    \item Step 3: 45-60 minutes (469 scenarios)
    \item Steps 4-5: 2-5 minutes
\end{itemize}

\end{frame}

% ============================================
% SECTION 15: LESSONS LEARNED
% ============================================

\begin{frame}{Lessons Learned \& Best Practices}
\textbf{Key Insights from Development:}

\vspace{0.3cm}

\begin{enumerate}
    \item \textbf{Validate input data early}
    \begin{itemize}
        \item Discovered UGA/LBY issues mid-analysis
        \item Cost: Had to re-run and document
        \item Lesson: Check GLORIA data quality first
    \end{itemize}

    \vspace{0.2cm}

    \item \textbf{Document as you go}
    \begin{itemize}
        \item Created WORK\_LOG.md during development
        \item Tracked all decisions and bug fixes
        \item Essential for replication
    \end{itemize}

    \vspace{0.2cm}

    \item \textbf{Test with single scenario first}
    \begin{itemize}
        \item Built \texttt{TEST\_ONE\_SCENARIO.py}
        \item Validated logic before batch processing
        \item Saved hours of debugging
    \end{itemize}

    \vspace{0.2cm}

    \item \textbf{Preserve base MINDSET modules}
    \begin{itemize}
        \item Minimal modifications to core files
        \item Custom logic in separate batch script
        \item Easier maintenance and updates
    \end{itemize}
\end{enumerate}

\end{frame}

% ============================================
% SECTION 16: FUTURE EXTENSIONS
% ============================================

\begin{frame}{Potential Extensions to This Methodology}
\textbf{Immediate Extensions:}

\vspace{0.3cm}

\begin{itemize}
    \item \textbf{Proxy country implementation}
    \begin{itemize}
        \item Replace UGA/LBY coefficients with Kenya/Tunisia
        \item Re-run batch with corrected data
        \item Sensitivity analysis comparing original vs proxy
    \end{itemize}

    \vspace{0.2cm}

    \item \textbf{Additional outcome variables}
    \begin{itemize}
        \item Value-added by sector
        \item CO$_2$ emissions (if GLORIA env extensions available)
        \item Import requirements
    \end{itemize}
\end{itemize}

\vspace{0.5cm}

\textbf{Methodological Extensions:}

\vspace{0.3cm}

\begin{itemize}
    \item \textbf{Dynamic analysis}
    \begin{itemize}
        \item Multi-year impacts (requires time-series MRIO)
        \item Capital depreciation effects
    \end{itemize}

    \item \textbf{Uncertainty quantification}
    \begin{itemize}
        \item Monte Carlo on employment coefficients
        \item Confidence intervals on job estimates
    \end{itemize}
\end{itemize}

\end{frame}

% ============================================
% SECTION 17: SUMMARY
% ============================================

\begin{frame}{Summary: What We Built}
\textbf{Four Major Methodological Contributions:}

\vspace{0.3cm}

\begin{enumerate}
    \item \textbf{Product Allocation System}
    \begin{itemize}
        \item Eora26 sectors → GLORIA products
        \item Country-specific distributions
    \end{itemize}

    \vspace{0.2cm}

    \item \textbf{"Strategy as Purchases" Implementation}
    \begin{itemize}
        \item Correct investment interpretation
        \item Uses MINDSET investment converter properly
    \end{itemize}

    \vspace{0.2cm}

    \item \textbf{Employment-Only Batch Processing}
    \begin{itemize}
        \item Streamlined workflow (95\% time reduction)
        \item Direct vs indirect separation
    \end{itemize}

    \vspace{0.2cm}

    \item \textbf{ISIC Aggregation Pipeline}
    \begin{itemize}
        \item Two-stage mapping to standard classifications
        \item Validation checks built in
    \end{itemize}
\end{enumerate}

\vspace{0.5cm}

\textbf{Output:} Fully replicable workflow documented with:
\begin{itemize}
    \item Custom scripts with line-by-line comments
    \item WORK\_LOG.md tracking all decisions
    \item Validation checks at each step
\end{itemize}

\end{frame}

% ============================================
% SECTION 18: DOCUMENTATION FILES
% ============================================

\begin{frame}{Complete Documentation Package}
\textbf{Location:} \texttt{Claude Code/temp/}

\vspace{0.3cm}

\textbf{Technical Documentation:}
\begin{enumerate}
    \item \texttt{MINDSET\_Workflow\_Guide.md}
    \begin{itemize}
        \item Step-by-step methodology
        \item Code references with line numbers
        \item Uganda/Libya anomaly analysis
    \end{itemize}

    \item \texttt{WORK\_LOG.md}
    \begin{itemize}
        \item Session-by-session development log
        \item Bug fixes and solutions
        \item Test results
    \end{itemize}
\end{enumerate}

\vspace{0.3cm}

\textbf{Scripts:}
\begin{enumerate}
    \setcounter{enumi}{2}
    \item \texttt{sector\_to\_product\_allocation.Rmd}
    \item \texttt{RunMINDSET\_EmploymentOnly\_BATCH\_FINAL.py}
    \item \texttt{CALCULATE\_L\_BASE\_ONCE.py}
    \item \texttt{TEST\_ONE\_SCENARIO.py} (for validation)
\end{enumerate}

\vspace{0.3cm}

\textbf{Reference Documents:}
\begin{enumerate}
    \setcounter{enumi}{6}
    \item \texttt{Uganda\_Libya\_Anomaly\_Summary.md}
    \item This presentation
\end{enumerate}

\end{frame}

% ============================================
% CLOSING
% ============================================

\begin{frame}[plain]
\centering
\vspace{1.5cm}
{\Huge Questions?}

\vspace{1.5cm}

\textbf{Key Takeaway:}

\vspace{0.3cm}

We built a custom employment-only mode for MINDSET that:
\begin{itemize}
    \item Processes 469 scenarios in $<$1 hour (vs 78 hours)
    \item Correctly interprets investment as product demand
    \item Maps results to standard ISIC classifications
    \item Fully documented and replicable
\end{itemize}

\vspace{1cm}

\small
Contact: Felipe Esteves \\
RAND Corporation \\
Documentation: \texttt{Claude Code/temp/}

\end{frame}

\end{document}
